\documentclass[12pt,a4paper]{article}

\usepackage[utf8]{inputenc}
\usepackage[T1]{fontenc}
\usepackage[margin=2.5cm]{geometry}
\usepackage{hyperref}
\usepackage{microtype}
\usepackage{parskip}
\usepackage{lmodern}

\hypersetup{colorlinks=true, urlcolor=blue!60!black}

\pagestyle{empty}

\begin{document}

\begin{flushright}
Kundai Farai Sachikonye \\
Auf der Lichtung 19, 82178 Puchheim, Germany \\
\href{https://github.com/fullscreen-triangle}{github.com/fullscreen-triangle} \\
\texttt{kundai.sachikonye@bitspark.com} \\
(+49) 1626386344 \\[6pt]
\today
\end{flushright}

\vspace{1em}

AMBIOM Research Group \\
Leibniz-Institut für Analytische Wissenschaften --- ISAS --- e.V. \\
Dortmund

\vspace{1.5em}

\textbf{Re: PhD Candidate, AMBIOM (Analysis of Microscopic BIOMedical images) ---
Reference 401\_2026}

\vspace{1em}

Dear Hiring Team,

I am applying for the PhD Candidate position in the AMBIOM group. The three tasks
in the posting --- a large-scale microscopy image-analysis pipeline, a
computational framework for 3D shape modelling, and AI-based representation
learning for cell modelling --- describe almost exactly the work I have been
doing, and I would be glad to develop it into a doctorate at ISAS.

\textbf{Microscopy image analysis and computer vision.}
This is my core. I have built a quantitative microscopy image-analysis framework
covering segmentation, point-spread-function deconvolution, registration, and
morphometry, with uncertainty quantification, for fluorescence and phase-contrast
microscopy. It is GPU-accelerated (WebGL/GLSL) and deployed as a live tool
(\href{https://hieronymus-omega.vercel.app/}{hieronymus-omega.vercel.app}); the
underlying method is written up as a formal image-analysis paper. Extending this
pipeline to large-scale 3D microscopy is a direct continuation of what I already
do.

\textbf{Shape modelling and representation learning.}
My image-analysis work centres on recovering metric-preserving geometry from
image observations --- scale fields, coordinate fields, geodesic distance --- which
is the natural basis for quantitative 3D shape characterisation. On the learning
side, I work in PyTorch with CNNs, autoencoders and representation learning, and
foundation-model fine-tuning, applied to imaging and spectral data and to
data-driven cell- and disease-state modelling. Building a representation-learning
framework for cell modelling is squarely in this line.

\textbf{Biology and communication.}
I hold an MSc in Computational Biology and an MSc in Pharmaceutical Biotechnology,
with a BSc in Biochemistry and Cell Biology, so the biological context of cell
modelling is familiar ground rather than a hurdle. I write and publish regularly,
and I am used to presenting technical work to interdisciplinary audiences. My
degree qualifies me directly (MSc in computational biology); English is native;
and I am based in Germany.

I am proactive and independent by temperament --- most of the work above was
self-directed --- and I would value the structure, supervision, and
state-of-the-art imaging infrastructure that a PhD at ISAS offers to turn it into
rigorous, publishable research. I would welcome the opportunity to discuss the
position.

\vspace{1.5em}
Yours sincerely,

\vspace{2.5em}
\textbf{Kundai Farai Sachikonye}

\end{document}
